\documentclass[11pt]{article}
\usepackage[margin=1in]{geometry}
\usepackage{amsmath,amssymb,booktabs}
\usepackage{natbib}
\usepackage[hidelinks]{hyperref}

\title{Decision Debt: A Number for Silent Decision Change Across Model Upgrades}
\author{LaxmiKumar Sammeta}
\date{2026}

\begin{document}
\maketitle

\begin{abstract}
An upgrade is judged by the headline number: accuracy on a fixed suite. We
built the number that accuracy hides: \textbf{decision debt}, how much the
model changes decisions, and how much of that change is a regression.

We route 300 cases from the BANKING77 intent set to six actions. Each case
carries an expected cost if the wrong action is chosen. We run five models
locally, on two routing prompts, and compare them on accuracy and on
decision change.

Two results:

\begin{enumerate}
\item \texttt{qwen2.5:0.5b} to \texttt{qwen2.5:1.5b}: accuracy \textbf{30.3\% to 61.0\%}
  (a 30.7-point gain) while \textbf{70.0\%} of decisions change.
\item \texttt{qwen2.5:1.5b} to \texttt{qwen2.5:3b}: accuracy \textbf{61.0\% to 59.3\%}
  (not significant, $p=0.50$) while \textbf{19.7\%} of decisions still change,
  with 6.7\% regressions and 5.0\% improvements.
\end{enumerate}

For the second model pair, the interesting cases are the ones where the
changed decisions go the wrong way. A team comparing those two models would
report an upgrade that changed nothing. We report churn and regressions as
first-class numbers, weighted by the cost they cause, per 1,000 decisions.
A fraud-detection agent built the same way is scored separately on the suite.
\end{abstract}

\section{Why accuracy is the wrong summary for an upgrade}

A model turns an input into an action. The same input, run through two
models, gives a verdict of ``same'' or ``different.'' An upgrade changes a
decision when that verdict is ``different.'' Accuracy only looks at
``correct'' versus ``wrong,'' and it is the average of the two: regressions
cancel improvements. Four quantities matter, and only one of them is in every
report.

\begin{enumerate}
\item \textbf{Churn.} The share of decisions that change between the two
  models. A model that changes 15\% of decisions is not equivalent, even if
  the average score is the same.
\item \textbf{Regressions.} The share that change from correct to wrong.
  Churn without regressions is a pure gain; churn with regressions is
  quality that moves in both directions.
\item \textbf{Cost-weighted churn.} The expected cost carried by the changed
  decisions, divided by the total cost exposure of the suite, normalized.
  Some actions cost more than others; see the cost model below.
\item \textbf{Decision loss.} The difference of expected cost per 1,000
  decisions between the two models. Signed: positive means the new model
  costs more per decision on average.
\end{enumerate}

The MCNemar test applies to the accuracy side, not the decision side. On the
upgrades below, it reports whether the accuracy change is distinguishable
from noise. Decisions are flows, not states, so the flow rates need their
own statistics: Wilson intervals on churn and regressions.

\section{Related work}

Upgrades are usually read through aggregate benchmarks. A few works have
turned to item-level and cost-weighted views.

\begin{itemize}
\item \textbf{Item-level regressions in API migrations.}
  \citep{api_migration_regressions} measures per-item changes across
  commercial API upgrades, classifying items as reliably improved, reliably
  regressed, or practically equivalent. Our unit is the business decision: a
  decision changes when the predicted action differs between two deployed
  models.
\item \textbf{Migration benchmarks.} \citep{upgradebench} is a
  decision-centric benchmark over migration choices for fine-tuned
  specialists; it measures migration decisions, not decision change between
  deployed models.
\item \textbf{Cost-aware abstention.} \citep{cost_aware_abstention} measures
  when to defer in IT ticket classification; the cost model we use measures
  the cost of decisions that change.
\item \textbf{Decision-model benchmarks.} \citep{general_decision_models}
  and \citep{decision_models_triage} benchmark decision models on triage
  accuracy, calibration, and out-of-scope detection.
\end{itemize}

None of these defines decision debt: the churn of decisions between two
versions, the regressions, and a cost-weighted version, with an operational
upgrade-gate framing.

\section{The cost model}

Routing a case to the wrong action is not free, and not equally bad on every
case. One formula and a table parameterize both sides.

\begin{equation}
\label{eq:cost}
\text{cost}(t, p) =
\begin{cases}
0 & \text{if } p = t,\\
\text{delay}[t] + \text{handling}[p] & \text{otherwise},
\end{cases}
\end{equation}

where $t$ is the gold action and $p$ the predicted action. \texttt{delay} is
what it costs to sit in the wrong team; \texttt{handling} is what that team
pays to touch the case.

\begin{table}[h]
\centering
\begin{tabular}{lrr}
\toprule
action $t$ & delay & handling \\
\midrule
\texttt{fraud} & 60 & 30 \\
\texttt{disputes} & 25 & 20 \\
\texttt{payments} & 18 & 14 \\
\texttt{cards} & 12 & 12 \\
\texttt{fees} & 7 & 9 \\
\texttt{account} & 5 & 8 \\
\bottomrule
\end{tabular}
\caption{Cost model parameters. \texttt{fraud} to \texttt{account} costs 68;
\texttt{account} to \texttt{fraud} costs 35. The matrix is not symmetric, so
``wrong'' is not one number.}
\end{table}

The asymmetry is not an option. A disputed charge that the wrong team sends
to account support costs a human to review it; an account-support case
forwarded to fraud costs the same team's time. We take the difference as
data and the numbers as a parameterization, not as ground truth.

\section{The suite}

Source: the BANKING77 train split of PolyAI-LDN/task-specific-datasets, 10,003
customer messages over 77 intents.

\begin{enumerate}
\item The 77 intents map to six actions: \texttt{fraud}, \texttt{disputes},
  \texttt{payments}, \texttt{cards}, \texttt{fees}, \texttt{account}. The
  mapping table and the rationale are in the code.
\item 300 cases, 50 per action, sampled with a fixed seed.
\item Each case carries: the message text, the gold action, and a six-entry
  cost vector following \texttt{cost} in \eqref{eq:cost}.
\end{enumerate}

\section{Models and method}

Five models, all local, all through Ollama on the local machine (Apple M4,
16 GB RAM):

\begin{tabular}{lrr}
model & size \\
\hline
\texttt{llama3.1:8b} & 4.9 GB \\
\texttt{qwen2.5-coder:7b} & 4.7 GB \\
\texttt{qwen2.5:0.5b} & 0.4 GB \\
\texttt{qwen2.5:1.5b} & 1.0 GB \\
\texttt{qwen2.5:3b} & 1.9 GB \\
\end{tabular}

Two routing prompts. \texttt{v1} asks the model to answer with one of the six
actions. \texttt{v2} adds two rules (choose \texttt{fraud} for missing money
or a possibly compromised card; choose \texttt{disputes} for money-back
requests) and demands JSON output. Prompt changes are not part of
the story here: the relevant comparison is the same model with the two
prompts, to separate prompt sensitivity from model sensitivity.

Decoding: temperature 0, seed 42, the output constrained to the six allowed
labels. Of the 2,100 decision calls, five returned no parseable label
(answers like ``disposables'' or ``balances'') and were counted as wrong.

\section{Runs}

\begin{table}[h]
\centering
\begin{tabular}{lrr}
\toprule
run & accuracy & \$/1,000 \\
\midrule
\texttt{llama3.1:8b} v1 & 66.0\% & 14,473 \\
\texttt{llama3.1:8b} v2 & 66.3\% & 14,480 \\
\texttt{qwen2.5:0.5b} v1 & 30.3\% & 27,607 \\
\texttt{qwen2.5:1.5b} v1 & 61.0\% & 16,643 \\
\texttt{qwen2.5:1.5b} v2 & 62.3\% & 16,683 \\
\texttt{qwen2.5:3b} v1 & 59.3\% & 17,500 \\
\texttt{qwen2.5-coder:7b} v1 & 67.0\% & 13,977 \\
\bottomrule
\end{tabular}
\caption{Accuracy and expected cost per 1,000 decisions per run.}
\end{table}

\section{Upgrades}

\begin{table}[h]
\centering
\begin{tabular}{llllll}
\toprule
pair & acc $\Delta$ & churn & regressions & cw churn & $p$ \\
\midrule
llama3.1:8b v1→v2 (prompt) & +0.3pp & 12.3\% [9.1, 16.5] & 4.7\% & 6.2\% & 1.00 \\
llama3.1:8b → 0.5b & $-35.7$pp & 71.7\% [66.3, 76.5] & 44.0\% & 47.6\% & $\sim 0$ \\
llama3.1:8b → 1.5b & $-5.0$pp & 26.7\% [22.0, 31.9] & 11.0\% & 17.8\% & 0.049 \\
1.5b → 3b & $-1.7$pp & 19.7\% [15.6, 24.5] & 6.7\% & 12.5\% & 0.50 \\
llama3.1:8b → coder & +1.0pp & 18.7\% [14.7, 23.5] & 5.7\% & 10.1\% & 0.74 \\
\bottomrule
\end{tabular}
\caption{Upgrade pairs: accuracy change, churn, regressions, cw churn.}
\end{table}

What to read in this table.

The first row is the control: the same model, a new prompt, no accuracy
change, 12.3\% of decisions changed, and the accuracy test says it is noise
($p = 1.00$). Churn is present.

\texttt{qwen2.5:1.5b} to \texttt{qwen2.5:3b} is the quiet one: accuracy moves
by less than 2 points and is not significant ($p = 0.50$), while 19.7\% of
decisions change, with 6.7\% regressions and 5.0\% improvements. A team
comparing these two models by accuracy would see no upgrade.

\texttt{qwen2.5:0.5b} to \texttt{qwen2.5:1.5b} is the loud one: accuracy rises
by 30.7 points and 70.0\% of decisions change; 8.7\% of decisions go from
correct to wrong and 39.3\% become correct where they were wrong. The ledger
is \texttt{qwen2.5:0.5b} at \$27,607 per 1,000 decisions,
\texttt{qwen2.5:1.5b} at \$16,643.

\section{The fraud-detection agent}

The same suite, asked a different question: build a local fraud-detection
agent that answers \texttt{fraud} or \texttt{not\_fraud} per message, and score
it on the suite. Gold is the suite's gold action (50 fraud cases of 300);
the raw runs are in \texttt{out/fraud/}.

\begin{table}[h]
\centering
\begin{tabular}{lrrr}
\toprule
model & overall acc (300) & fraud-case acc (50) & parse fails \\
\midrule
\texttt{llama3.1:8b} & 85.3\% & 32.0\% & 0 \\
\texttt{qwen2.5:0.5b} & 29.7\% & 74.0\% & 0 \\
\texttt{qwen2.5:1.5b} & 80.7\% & 52.0\% & 0 \\
\texttt{qwen2.5:3b} & 84.0\% & 6.0\% & 0 \\
\texttt{qwen2.5-coder:7b} & 89.0\% & 54.0\% & 0 \\
\bottomrule
\end{tabular}
\caption{Fraud-agent accuracy: overall over all 300 cases, and of the 50
fraud cases.}
\end{table}

Stability of the fraud answer across models, churn against
\texttt{llama3.1:8b}:

\begin{table}[h]
\centering
\begin{tabular}{ll}
\toprule
model & churn of the fraud answer \\
\midrule
\texttt{qwen2.5:0.5b} & 71.7\% \\
\texttt{qwen2.5:1.5b} & 15.3\% \\
\texttt{qwen2.5:3b} & 7.3\% \\
\texttt{qwen2.5-coder:7b} & 9.0\% \\
\bottomrule
\end{tabular}
\end{table}

Reading: the fraud signal is a low-volume one. Of the 50 fraud cases,
\texttt{llama3.1:8b} catches 16 (32\%), \texttt{qwen2.5-coder:7b} catches 27
(54\%), \texttt{qwen2.5:1.5b} catches 26 (52\%), and \texttt{qwen2.5:0.5b}
catches 37 (74\%). The smallest model catches the most fraud, though its
overall accuracy is only 29.7\%. Against \texttt{llama3.1:8b}, the fraud
answer churns in 71.7\% of cases from \texttt{qwen2.5:0.5b} and 15.3\% from
\texttt{qwen2.5:1.5b}: a model can be a better fraud detector and a less
stable fraud source.

\section{Code and data}

The paper is reproduced from
\url{https://github.com/slkreddy/decision-debt}.

The suite, the runs, and the analysis live in \texttt{out/}; the harness in
\texttt{src/}: \texttt{build\_cases.py}, \texttt{run\_eval.py},
\texttt{fraud\_eval.py}, \texttt{analyze.py}.

\section{Limitations}

The suite is intent routing, not enterprise policy. The actions and the cost
table are a stated assumption, not measured ground truth. 300 cases is small;
the accuracy changes with the largest $p$ are the prompt change ($p = 1.00$),
the comparison to \texttt{qwen2.5-coder:7b} ($p = 0.74$), and the 1.5B to 3B
upgrade ($p = 0.50$); where churn is involved, conclusions should be tentative.

The fraud study uses the same 300-case suite; its labels are the gold
actions (50 fraud of 300) and it does not use the cost model, since it
measures detection accuracy rather than decision damage.

Churn and regressions are the statistically cleaner outcomes: they do not
depend on the gold labels beyond correctness. They are what you want for a
gate before a model goes into service.

\section{Conclusion}

The metric is decision debt, the combination of:

\begin{enumerate}
\item of all decisions, how many changed ($p$ churn);
\item of those, how many went the wrong way ($p$ regressions);
\item how much the moved decisions cost, relative to the suite's total cost
  exposure;
\item the dollar change per 1,000 decisions.
\end{enumerate}

Report them together. One number, ``did the upgrade change decisions?,''
replaces accuracy when the question is whether to ship. The four numbers are
the decision debt of the upgrade.

\begin{thebibliography}{9}
\bibitem[Xu and Wu, 2026]{api_migration_regressions}
Xiaonan Xu and Wenjing Wu. What aggregate scores miss: measuring item-level
regressions in commercial LLM API migrations.
\url{https://www.alphaxiv.org/abs/2608.17719}.

\bibitem[Chen and Zhang, 2026]{upgradebench}
Ye Chen and Weining Zhang. UpgradeBench: A decision-centric benchmark for
upgrading fine-tuned LLM specialists.
\url{https://www.alphaxiv.org/abs/2608.20918}.

\bibitem[Gupta, 2026]{cost_aware_abstention}
Sachin Gupta. Cost-aware abstention for LLM-based IT ticket classification.
\url{https://www.alphaxiv.org/abs/2610.31657}.

\bibitem[Duan et al., 2026]{general_decision_models}
Feiyu Duan, Jiayu Lin, Jia Wang, Jun Xiang, Jialiang Wu, Xinnong Zhang,
Hanqi Yan, Siyuan Wang, and Zhongyu Wei. General decision models:
benchmarking and insights beyond Jev.
\url{https://www.alphaxiv.org/abs/2610.03935}.

\bibitem[Roumeliotis et al., 2026]{decision_models_triage}
Konstantinos I. Roumeliotis, Dionisis Margaris, Dimitris Spiliotopoulos, and
Nikolaos D. Tselikas. Decision models for business customer-service triage:
accuracy, calibration, out-of-scope detection and cost of Jev, Clef and
Perplexity Decider.
\url{https://www.alphaxiv.org/abs/2611.customer-service-triage-decision-models}.

\bibitem[Banking77, 2026]{banking77}
PolyAI-LDN. task-specific-datasets: BANKING77 train split.
\url{https://github.com/PolyAI-LDN/task-specific-datasets}.
\end{thebibliography}

\end{document}
